% Generated by Listas/make_statements.py from pratica-lecture-06.md, item 1. Do not edit.
\textbf{Source.} Muñoz-Garcia \& Toro-Gonzalez (2019), \emph{Strategy and Game Theory: Practice Exercises with Answers}, Ch. 5, Exercise 4 ``Cartel with Firms Competing a la Cournot---Unsustainable in a One-Shot game'', p.~228 {[}PDF 240{]}. Reproduced verbatim (statement only; the book's worked answer follows it).



\begin{quote}

Consider two firms competing in output levels (a la Cournot) and facing linear inverse demand \(p(Q) = 100 - Q\), where \(Q = q_1 + q_2\) denotes aggregate output. For simplicity, assume that firms face a common marginal cost of production \(c = 10\).



\begin{enumerate}

\def\labelenumi{(\alph{enumi})}

\item

  \emph{Best response function.} Set up each firm's profit-maximization problem and find its best response function. Interpret.

\item

  \emph{Cournot competition.} Find the equilibrium output each firm produces when competing a la Cournot (that is, when they simultaneously and independently choose their output levels). In addition, find the profits that each firm earns in equilibrium.

\item

  \emph{Collusive agreement.} Assume now that the CEOs from both companies meet for lunch and start talking about how they could increase their profits if they could increase market prices by reducing their production. That is, firms are trying to collude! Set up the maximization problem that firms solve when maximizing their \emph{joint} profits (that is, the sum of profits for both firms). Find the output level that each firm should select to maximize joint profits. In addition, find the profits that each firm obtains in this collusive agreement.

\item

  \emph{Profit comparison.} Compare the profits that firms obtain when competing a la Cournot (from part b) against their profits when they successfully collude (from part c).

\item

  \emph{Unsustainable cartel.} Let us now show that the collusive agreement from part (c) cannot be sustained as a Nash equilibrium of the one-shot game, that is, when firms interact only once. Using the best response function found in part (a), show that when firm \(i\) chooses the collusive output found in part (c) its rival (firm \(j\)) responds with an output that is different from its collusive output. Interpret.

\end{enumerate}

\end{quote}

